\section{一般Hesse型の全混合方向に対する普遍中心三次式}
\label{sec:pg-hesse-mixed-universal}
本節ではFermat型に限らず、Hesse三次形式
\[
 g_t(x,y,z)=\frac{x^3+y^3+z^3}{3}+txyz
\]
の任意のパラメータ$t$に対して、中心変数を含む全ての三次混合項を
\emph{有限個の多項式係数}で量子化した中心三次元を与える。
この構成は形式的中心元の存在補題
\ref{lem:pg-formal-center}を使用しない。
非零の混合係数を独立な記号のまま保持し、
$\hbar,t$の有理関数体上で自由語の関係式への所属を
厳密に検証することで証明する。
得られる式はFermat型の
定理\ref{thm:pg-fermat-mixed-universal}を$t=0$として含む。

\subsection{短い多項式による普遍公式}
$R$を標数零の可換次数付き$\CC$-代数とする。
$A,B,C\in R_1$、$D,E,H\in R_2$、$G\in R_3$、
$t\in\CC$を任意に取って
\begin{equation}
 F=g_t(x,y,z)+Axy+Byz+Czx+Dx+Ey+Hz+G
 \label{eq:pg-hm-F}
\end{equation}
とおく。$R[\hbar]$上の次数付き代数
$\mathcal A_{F,\hbar}$を
\begin{equation}
\begin{aligned}
 [x,y]&=\hbar\bigl(z^2+t\sym(xy)+By+Cx+H\bigr),\\
 [y,z]&=\hbar\bigl(x^2+t\sym(yz)+Ay+Cz+D\bigr),\\
 [z,x]&=\hbar\bigl(y^2+t\sym(zx)+Ax+Bz+E\bigr)
\end{aligned}
\tag{HR}\label{eq:pg-hm-rel}
\end{equation}
で定める。$\deg\hbar=0$であり、全ての$R$の元は中心とする。
記法
\[
\begin{aligned}
 S={}&x^3+y^3+z^3,&
 T={}&\sym(xyz),\\
 Q={}&A\sym(xy)+B\sym(yz)+C\sym(zx),&
 Z_2={}&Az^2+Bx^2+Cy^2,\\
 L={}&Dx+Ey+Hz,&
 R_2={}&B^2x+C^2y+A^2z,\\
 R_{11}={}&ACx+ABy+BCz
\end{aligned}
\]
を用いる。ここで$\sym(xyz)$は六語の平均である。

\begin{theorem}[Hesse型の普遍混合中心三次式]
\label{thm:pg-universal-hesse-mixed}
$q=\hbar^2$とし、次の七つの多項式を定義する。
\begin{equation}
\begin{aligned}
 U={}&1+\frac{qt^2}{2}+\frac{q^2t}{6}-\frac{q^2t^4}{48},\\
 W={}&\frac{(4-qt^2)(12t+q(t^3-8))}{48},\\
 c={}&\frac{(4-qt^2)(12+qt^2)}{48},&
 d={}&\frac{q(12t+q(2-t^3))}{36},\\
 e={}&1+\frac{q^2t(8-t^3)}{144},&
 f={}&\frac{q(4-qt^2)}{24},\\
 g={}&-\frac{q(12t+q(t^3+4))}{72}.
\end{aligned}
\tag{HC}\label{eq:pg-hm-coefficients}
\end{equation}
このとき
\begin{equation}
 \boxed{\quad
 \Omega_{F,\hbar}
 =\frac{U}{3}S+WT+cQ+dZ_2+eL+fR_2+gR_{11}+G
 \quad}
 \label{eq:pg-hm-center}
\end{equation}
は$\hbar=0$を含む適切なZariski開部分上で
$\mathcal A_{F,\hbar}$の\emph{厳密な中心元}となる。
また$\Omega_{F,0}=F$である。
これは$\hbar$について有限な多項式公式であり、
$n$個の中心変数に係数$A,\ldots,G$を特殊化する場合でも
中心変数の個数に依存せず同じ表示が有効である。
$t^3=-1$を排除する条件はこの\emph{中心性の恒等式そのもの}には不要だが、
$g_t=0$が滑らかな平面三次であるためには$t^3\ne-1$が必要である。
\end{theorem}
\begin{proof}
基底係数を
$R_0=\mathbb Q[A,B,C,D,E,H]$、
$\deg A=\deg B=\deg C=1$、
$\deg D=\deg E=\deg H=2$とする。
$t,\hbar$は独立な次数零の不定元とする。
$G$は中心なので零としてよい。
非可換自由語では中央係数を可換な単項式としてまとめ、
$x,y,z$の順序付き語をそのまま保持する。
以下、全次数四の自由語空間を$V_4$、式
\eqref{eq:pg-hm-rel}の三つの二次関係式の
左右からの全次数二の倍数が張る空間を$I_4$とする。

中央係数の重み$0,1,2,3,4$の単項式の個数は
順に$1,3,9,19,39$である。各重み$c$について
$x,y,z$からなる長さ$4-c$の語を列挙すると
\[
 \dim_{\mathbb Q(\hbar,t)}V_4
 =1\cdot81+3\cdot27+9\cdot9+19\cdot3+39=339.
\]
全次数二の左右の倍数を列挙すると
162個の関係ベクトルが得られる。
\emph{これらは全次数四の二次関係式のイデアル成分を生成する。}

式\eqref{eq:pg-hm-coefficients}を
\eqref{eq:pg-hm-center}に代入し、三交換子
$[\Omega_{F,\hbar},x]$、$[\Omega_{F,\hbar},y]$、
$[\Omega_{F,\hbar},z]$を自由語で展開する。
162関係ベクトルの行列を$\mathbb Q(\hbar,t)$上で
厳密に行簡約すると
\[
 \operatorname{rank}I_4=144,\qquad
 [\Omega_{F,\hbar},x],\
 [\Omega_{F,\hbar},y],\
 [\Omega_{F,\hbar},z]\in I_4
 \tag{UC}\label{eq:pg-hm-certificate}
\]
となる。この恒等式は六つの中心係数を独立不定元として扱い、
$\hbar,t$を特殊化せずに得られた。
検証コード
\texttt{computations/verify\_universal\_hesse\_mixed\_center.py}
が全339基底、162関係と三交換子を再生成して
有理関数体上の行簡約を実行する。
有限体や有理数標本での補間は証明に用いない。

分数体上での関係式所属を$\hbar=0$まで降下させるには
階数を比較する。$\hbar=0$で二次関係は通常の可換化関係なので
\[
 \dim(V_4/I_4)|_{\hbar=0}
 =15+3\cdot10+9\cdot6+19\cdot3+39=195.
\]
従って$\operatorname{rank}I_4|_{\hbar=0}=339-195=144$であり、
上の一般階数と一致する。
ここで最後の等式は任意の$t$について成立する。
従って$(\hbar,t)$-平面の直線$\hbar=0$を含む
あるZariski開近傍で、関係行列の階数は144に一定であり、
階数144の小行列は単元として取れる。
行空間とその商はこの近傍上局所自由であるから、
分数体上で零の三交換子はその局所自由商でも零となる。
ゆえに任意の$t$に特殊化しても、
$\hbar=0$の適切な近傍で厳密な中心性が成立する。
この降下に、全次数のPBW性を別途仮定する必要はない。

最後に$q=0$で$U=c=e=1$、$W=t$、
$d=f=g=0$なので$\Omega_{F,0}=F$。
\end{proof}

\begin{corollary}[Hessianによるコンパクトな普遍式]
\label{cor:pg-hesse-hessian-form}
定理\ref{thm:pg-universal-hesse-mixed}の$F$に対し
\[
 H_F=\det\operatorname{Hess}_{x,y,z}(F)
\]
とおく。また上記の$Q,Z_2,L,R_2,R_{11}$を使い
\[
\begin{aligned}
 \Psi_t(F)={}&
 \frac{t(8-t^3)}{48}g_t
 -\frac{t^4}{48}Q
 +\frac{2-t^3}{36}Z_2\\
 &+\frac{t(8-t^3)}{144}L
 -\frac{t^2}{24}R_2
 -\frac{t^3+4}{72}R_{11}.
\end{aligned}
\]
すると普遍中心三次元は
\begin{equation}
 \boxed{\quad
 \Omega_{F,\hbar}
 =\sym\left(
 F-\frac{\hbar^2}{12}H_F
 +\frac{\hbar^2}{6}ABC
 +\hbar^4\Psi_t(F)
 \right).
 \quad}
 \label{eq:pg-hesse-hessian-closed}
\end{equation}
特に$\hbar^2$項は普遍Hessian補正
（命題\ref{prop:pg-universal-casimir-jet}）であり、
中心係数$ABC/6$は中心元の規格化の自由度を反映する。
本式は$\hbar^4$項まで含めることで
\emph{全ての$\hbar$次数について厳密に中心}となる。
\end{corollary}
\begin{proof}
$H_F$をactive変数だけで微分して展開すると
\[
 H_F=-2t^2S+2(t^3+4)xyz+
 2t^2Q-4tZ_2-2R_2+2tR_{11}+2ABC
\]
である。
これを右辺に代入して$\hbar^2$と$\hbar^4$の係数を比較すると、
式\eqref{eq:pg-hm-center}と恒等的に一致する。
従って定理の中心性から従う。
\end{proof}

\begin{remark}[既存の中心元との整合性]
$t=0$では
$U=1$、$W=-2\hbar^2/3$、$c=e=1$、
$d=\hbar^4/18$、$f=\hbar^2/6$、
$g=-\hbar^4/18$となる。
ゆえに定理\ref{thm:pg-fermat-mixed-universal}の
全混合Fermat式を正確に回復する。
$A=B=C=D=E=H=G=0$では
定理\ref{thm:explicit-hesse-center}の
Sklyanin中心元を回復する。
それらは既知の入力・先行研究上の貢献と区別する。
\end{remark}

\subsection{対角二次混合と任意の滑らかな三変数三次式}

\begin{corollary}[全ての中心混合項を含むHesse型]
\label{cor:pg-universal-hesse-all-mixed}
中心次数一の$\ell_x,\ell_y,\ell_z$に対し
\[
 \widetilde F=g_t(x,y,z)
 +\ell_xx^2+\ell_yy^2+\ell_zz^2
 +Axy+Byz+Czx+Dx+Ey+Hz+G
\]
を考える。$X=x+\ell_x$、$Y=y+\ell_y$、
$Z=z+\ell_z$と置く。係数を
\begin{equation}
\begin{aligned}
 A'&=A-t\ell_z,& B'&=B-t\ell_x,& C'&=C-t\ell_y,\\
 D'&=D-\ell_x^2-A\ell_y-C\ell_z+t\ell_y\ell_z,\\
 E'&=E-\ell_y^2-A\ell_x-B\ell_z+t\ell_x\ell_z,\\
 H'&=H-\ell_z^2-B\ell_y-C\ell_x+t\ell_x\ell_y,\\
 G'&=G+\tfrac23(\ell_x^3+\ell_y^3+\ell_z^3)
 +A\ell_x\ell_y+B\ell_y\ell_z+C\ell_z\ell_x\\
 &\qquad-D\ell_x-E\ell_y-H\ell_z
 -t\ell_x\ell_y\ell_z
\end{aligned}
\tag{HT}\label{eq:pg-hm-translation}
\end{equation}
と定めると
\[
 \widetilde F
 =g_t(X,Y,Z)+A'XY+B'YZ+C'ZX+D'X+E'Y+H'Z+G'.
\]
したがって定理\ref{thm:pg-universal-hesse-mixed}の公式に
$A',B',C',D',E',H',G'$を代入して
$X=x+\ell_x$などへ引き戻すと、
\emph{任意の}Hesse型三次式と中心変数の混合項に対する
有限多項式の厳密な中心三次元が得られる。
\end{corollary}
\begin{proof}
表示式は三次式の展開による恒等式である。
$\ell_x,\ell_y,\ell_z$は中心元なので、
平行移動はJacobian量子化の交換子を保ち、
次数付き代数の同型を定める。
従って中心元を引き戻せる。
\end{proof}

\begin{corollary}[任意の滑らかな平面三次制限]\label{cor:pg-smooth-ternary-explicit}
$F(x,y,z,u)$を斉次三次式とし、
$g(x,y,z)=F(x,y,z,0)$が滑らかな平面三次曲線を定めるとする。
Hesse標準形\cite[Lemma~1]{ArtebaniDolgachev}により、
$M\in\mathrm{GL}_3(\CC)$、$\lambda\ne0$、
$t^3\ne-1$を選んで
$g(x)=\lambda\,g_t(Mx)$と書ける。
\[
 F'(X,u)=\lambda^{-1}F(M^{-1}X,u),
 \qquad \hbar'=\lambda\det(M)\,\hbar
\]
とおく。$F'$はその混合係数を
系\ref{cor:pg-universal-hesse-all-mixed}の形に一意に書ける。
同系によって構成した中心元$\Omega_{F',\hbar'}$から
\[
 \boxed{\quad
 \Omega_{F,\hbar}(x,u)
 =\lambda\Omega_{F',\hbar'}(Mx,u)
 \quad}
 \tag{HG}\label{eq:pg-hm-arbitrary-ternary}
\]
を得る。これは$F$のJacobian量子化において
$\hbar=0$近傍の\emph{厳密な}中心三次元であり、
$\Omega_{F,0}=F$である。
形式的中心元の存在補題を入力せず、
Hesse標準形の選択と有限個の多項式演算だけで構成できる。
\end{corollary}
\begin{proof}
定数三重ベクトルの変換則から、
$X=Mx$の座標ではJacobian交換子の量子化パラメータが
$\hbar'=\lambda\det(M)\hbar$になる。
中心元は代数同型で中心元に移り、
$\lambda$倍によって零ファイバーの剰余が$F$となる。
\end{proof}

\begin{corollary}[滑らかな制限を持つ全Poissonブロックの明示的代数的量子化]
\label{cor:pg-smooth-block-algebraic-quant}
定理\ref{thm:pg-normal}の正規形
$F=\sum_{s=1}^k f_s(x_s,u)$について、
各$g_s=f_s(x_s,0)$が滑らかな平面三次曲線を定めるとする。
各ブロックで系\ref{cor:pg-smooth-ternary-explicit}の
有限多項式中心元$\Omega_{s,\hbar}$を構成すれば
\[
 \Omega_\hbar=\sum_{s=1}^k\Omega_{s,\hbar},
 \qquad \Omega_0=F
\]
は代数的な厳密中心三次元となる。
命題\ref{prop:pg-PBW}の開集合を必要に応じて縮めると
$\Omega_\hbar$は各ファイバーで左右正則であり、
周囲の量子化代数はKoszul AS-regular、
中心商は定理\ref{thm:pg-quantization}の
Hilbert級数とAS-Gorenstein条件を満たす。
この部分族に対する量子化の構成は
形式的中心元存在補題を全く使用しない。
\end{corollary}
\begin{proof}
各ブロックに中心変数を共有させ、
異なるブロックのactive生成元を可換と取る。
個別の中心元は他ブロックと可換なので、その和は中心である。
各$g_s$が滑らかであるから、命題\ref{prop:pg-PBW}は
各ブロックのPBW自由性と$u=0$でのAS正則性を与える。
定理\ref{thm:pg-quantization}の証明のうち
中心元存在の補題の利用を
系\ref{cor:pg-smooth-ternary-explicit}に置き換える。
残る中央正則拡大、active次数filtrationによる
$\Omega_\hbar$の左右正則性、
中心三次商のAS-Gorenstein性の証明はそのまま適用できる。
\end{proof}

\begin{remark}[範囲と新規性]
本定理は、非可換量子化の中心性を\emph{全係数の恒等式}として
証明するものであって、各閉ファイバーのAS正則性・domain性や
中心商の平坦性をこの有限次数の検算のみから結論しない。
それらは命題\ref{prop:pg-PBW}、定理\ref{thm:pg-quantization}で
別途扱う。
一つの固定した元の座標に対するGL$_3$同変な短い不変式を
得たとは主張しない。座標選択でHesse標準形を求めた後に
適用できる\emph{有限多項式の統一的公式}である。
また$g=0$、節点型・尖点型など、三変数制限が滑らかでない
場合の包括的な閉公式は本節からは従わない。
公式自体の先行研究上の完全な新規性は未確認である。
\end{remark}
